% القسم: المناهج
% الفصل: الاستقراء
% المبحث: الاستقراء القوي

\documentclass[../../../include/open-logic-section]{subfiles}

\begin{document}

\olfileid{mth}{ind}{str}

\olsection{الاستقراء القوي}

كان مبدأ الاستقراء المتقدم يجمع بين إثبات $\OLId{latin-looped}{P}(\OLMathNumeral{0})$ وإثبات أن صدق
$\OLId{latin-looped}{P}(\OLId{latin-ordinary}{k})$ يستلزم صدق $\OLId{latin-looped}{P}(\OLId{latin-ordinary}{k}+\OLMathNumeral{1})$. وفي القسم الثاني يُفرض $\OLId{latin-looped}{P}(\OLId{latin-ordinary}{k})$ ويُتخذ
سبيلًا إلى إثبات $\OLId{latin-looped}{P}(\OLId{latin-ordinary}{k}+\OLMathNumeral{1})$. وهذا يكافئ أن نفرض، لكل $\OLId{latin-ordinary}{k}>\OLMathNumeral{0}$،~$\OLId{latin-looped}{P}(\OLId{latin-ordinary}{k}-\OLMathNumeral{1})$
فنثبت منه $\OLId{latin-looped}{P}(\OLId{latin-ordinary}{k})$. فالمعتبر إمكان الانتقال من كل عدد إلى خلفه، أي
إثبات القضية للعدد على فرض ثبوتها لسابقه.

وللمبدأ صورة أخرى لا نقتصر فيها على فرض القضية لسابق $\OLId{latin-ordinary}{k}$، أعني
$\OLId{latin-ordinary}{k}-\OLMathNumeral{1}$، بل نفرضها لجميع الأعداد الأصغر من~$\OLId{latin-ordinary}{k}$، ثم نستخرج منها القضية
لـ$\OLId{latin-ordinary}{k}$. وهذه الصورة تثبت كذلك $\OLId{latin-looped}{P}(\OLId{latin-ordinary}{n})$ لكل~$\OLId{latin-ordinary}{n} \in \Nat$. فإذا ثبت
$\OLId{latin-looped}{P}(\OLMathNumeral{0})$ فقد ثبت أن $\OLId{latin-looped}{P}$ تصدق لجميع الأعداد الأصغر من~$\OLMathNumeral{1}$. ثم إذا كان
معلومًا أن صدق $\OLId{latin-looped}{P}(\OLId{latin-ordinary}{l})$ لكل $\OLId{latin-ordinary}{l}<\OLId{latin-ordinary}{k}$ يوجب صدق $\OLId{latin-looped}{P}(\OLId{latin-ordinary}{k})$، شمل العلم الحالة
$\OLId{latin-ordinary}{k}=\OLMathNumeral{1}$؛ فنستنتج $\OLId{latin-looped}{P}(\OLMathNumeral{1})$. وبذلك ثبتت $\OLId{latin-looped}{P}(\OLMathNumeral{0})$ و$\OLId{latin-looped}{P}(\OLMathNumeral{1})$، أي ثبتت
$\OLId{latin-looped}{P}(\OLId{latin-ordinary}{l})$ لكل $\OLId{latin-ordinary}{l}<\OLMathNumeral{2}$. ومع الشرطية القائلة إن صدق $\OLId{latin-looped}{P}(\OLId{latin-ordinary}{l})$ لكل $\OLId{latin-ordinary}{l}<\OLMathNumeral{2}$ يوجب
صدق $\OLId{latin-looped}{P}(\OLMathNumeral{2})$، نستنتج~$\OLId{latin-looped}{P}(\OLMathNumeral{2})$، ثم نمضي على هذا النسق.

بل إذا برهنا الشرطية العامة «لكل $\OLId{latin-ordinary}{k}$، إذا صدقت $\OLId{latin-looped}{P}(\OLId{latin-ordinary}{l})$ لكل $\OLId{latin-ordinary}{l}<\OLId{latin-ordinary}{k}$
صدقت $\OLId{latin-looped}{P}(\OLId{latin-ordinary}{k})$»، استغنينا عن إثبات $\OLId{latin-looped}{P}(\OLMathNumeral{0})$ منفردة، لأنها تندرج فيها.
وذلك أن القول الكلي «لكل $\OLId{latin-ordinary}{l}<\OLId{latin-ordinary}{k}$،~$\OLId{latin-looped}{P}(\OLId{latin-ordinary}{l})$» صادق إذا لم يوجد أي $\OLId{latin-ordinary}{l}<\OLId{latin-ordinary}{k}$؛
وهذا هو التكميم الخالي. فقولنا «كل $\OLId{latin-looped}{A}$ هي $\OLId{latin-looped}{B}$» يصدق عند انتفاء كل
$\OLId{latin-looped}{A}$، وتصدق $\lforall[\OLId{latin-ordinary}{x}][(!A(\OLId{latin-ordinary}{x}) \lif !B(\OLId{latin-ordinary}{x}))]$ إذا لم يكن من $\OLId{latin-ordinary}{x}$ ما
يحقق~$!A(\OLId{latin-ordinary}{x})$. وصورة قولنا هنا هي
«$\lforall[\OLId{latin-ordinary}{l}][(\OLId{latin-ordinary}{l} < \OLId{latin-ordinary}{k} \lif \OLId{latin-looped}{P}(\OLId{latin-ordinary}{l}))]$»، وهي صادقة إذا خلا المجال من
$\OLId{latin-ordinary}{l} < \OLId{latin-ordinary}{k}$. وهذا واقع عند $\OLId{latin-ordinary}{k}=\OLMathNumeral{0}$؛ إذ لا يوجد $\OLId{latin-ordinary}{l}<\OLMathNumeral{0}$، فيصدق «لكل $\OLId{latin-ordinary}{l}<\OLMathNumeral{0}$،
$\OLId{latin-looped}{P}(\OLId{latin-ordinary}{l})$» مهما تكن $\OLId{latin-looped}{P}$. فبرهان «إذا كانت $\OLId{latin-looped}{P}(\OLId{latin-ordinary}{l})$ صادقة لكل $\OLId{latin-ordinary}{l}<\OLId{latin-ordinary}{k}$ فإن
$\OLId{latin-looped}{P}(\OLId{latin-ordinary}{k})$ صادقة» يتكفل تلقائيًّا بالحالة~$\OLId{latin-looped}{P}(\OLMathNumeral{0})$.

وتظهر فائدة هذه الصورة حين لا يكفي في إثبات القضية لـ$\OLId{latin-ordinary}{k}$ أن تُفرض
لـ$\OLId{latin-ordinary}{k}-\OLMathNumeral{1}$ وحده، بل يحتاج البرهان إلى صدقها لواحد أو أكثر من الأعداد
$\OLId{latin-ordinary}{l}<\OLId{latin-ordinary}{k}$.

\end{document}
